\chapter{Simplification and Subdivision}\label{ch:ops}

A mesh produced by a scanner or by marching cubes often has far more triangles than needed for
rendering, transmission or simulation; a mesh designed by an artist may have too few to look
smooth. \emph{Simplification} reduces the number of elements while preserving the shape as well
as possible, and \emph{subdivision} refines a coarse control mesh towards a smooth limit surface.
This chapter covers both, following Tema 8.

\section{Why piecewise linear surfaces}

In one dimension, approximating a function $f$ from samples at spacing $h$ by linear interpolation
has error $O(h^2)$, and interpolation by a polynomial of degree $p$ has error
$\frac{f^{(p+1)}(\xi)}{(p+1)!}\prod_i(x-x_i)=O(h^{p+1})$ by the generalised mean value theorem. Smooth
high-order representations are therefore more accurate per sample, which is the argument for
NURBS in CAD: they are inherently smooth and offer intuitive control through a lattice of control
points. However, a NURBS patch is limited to a grid-shaped domain and has the topology of a sheet,
a cylinder or a torus, so complex shapes need many patches with special rules to join them, and
deformations produce artefacts at the seams. Triangle meshes are inherently discrete, but need no
rules to join pieces, represent arbitrary topology, can be refined wherever resolution is needed,
can be made smooth by subdivision and are easy to render. Their approximation error is $O(h^2)$:
doubling the resolution in each direction (four times the vertices) reduces the error four times.
Triangles are the simplest linear element, are easy to reconstruct from point clouds and flexible
under deformation; quad meshes are preferred in animation and subdivision modelling but usually
require human intervention to build.

\section{Mesh simplification}

\subsection{Strategies and error metrics}

Three families of algorithms exist:
\begin{itemize}
  \item \textbf{Vertex clustering} \citep{Rossignac1993}: overlay a grid, merge all vertices in a
  cell into a representative, and remove degenerate triangles. It is fast, linear time and robust
  to any input, but the result depends on the grid, ignores the shape, and may change the
  topology.
  \item \textbf{Incremental decimation}: repeatedly remove the element whose removal causes the
  smallest error, by vertex removal and re-triangulation of the hole \citep{Schroeder1992} or,
  most commonly, by edge collapse.
  \item \textbf{Resampling}: compute a new mesh with the desired density (remeshing,
  \cref{ch:remesh}).
\end{itemize}
The quality of a simplified mesh $\M'$ is measured by its distance to the original $\M$, usually
the symmetric Hausdorff distance
$d_H(\M,\M')=\max\bigl(\sup_{\bx\in\M}d(\bx,\M'),\sup_{\bx\in\M'}d(\bx,\M)\bigr)$ or the mean
distance, estimated by dense sampling \citep{Cignoni1998}.

\subsection{Incremental decimation with edge collapses}

All incremental methods share the structure
\begin{enumerate}
  \item evaluate the local error of every removable element;
  \item while the target size is not reached: select the element with the smallest error, remove
  it, and update the error of the elements around it.
\end{enumerate}
With half-edge collapses ($\bx_1\to\bx_2$) the surviving vertex keeps its position; with full edge
collapses the merged vertex $\bar\bx$ may be placed anywhere, which allows better approximations.
Contracting edges in random order, without a geometric criterion, quickly destroys the shape; the
criterion is what matters.

\subsection{Quadric error metrics}

\Citet{Garland1997} measure the error of a vertex position as its squared distance to a set of
planes. A plane $ax+by+cz+d=0$ with unit normal $(a,b,c)$ is written $\bp=(a,b,c,d)\T$, and a point
in homogeneous coordinates $\bv=(x,y,z,1)\T$ has squared distance
\begin{equation}
(\bp\T\bv)^2=\bv\T(\bp\bp\T)\bv=\bv\T\mK_\bp\bv ,\qquad \mK_\bp=\bp\bp\T .
\end{equation}
Each vertex is associated with the planes of its incident triangles, and its \emph{quadric} is
the $4\times4$ symmetric matrix $\mQ_v=\sum_\bp\mK_\bp$ (optionally weighted by the triangle
areas). The error of placing the vertex at $\bv$ is $\Delta(\bv)=\bv\T\mQ_v\bv$, the sum of squared
distances to its planes; it is zero at the original position.

When an edge $(v_1,v_2)$ is collapsed into $\bar v$, the new vertex inherits the union of the planes,
approximated by $\bar\mQ=\mQ_1+\mQ_2$, and its position minimises
$\Delta(\bar\bv)=\bar\bv\T\bar\mQ\bar\bv$. The gradient condition
$\partial\Delta/\partial x=\partial\Delta/\partial y=\partial\Delta/\partial z=0$ is linear:
\begin{equation}
\begin{pmatrix}q_{11}&q_{12}&q_{13}&q_{14}\\q_{12}&q_{22}&q_{23}&q_{24}\\q_{13}&q_{23}&q_{33}&q_{34}\\0&0&0&1\end{pmatrix}\bar\bv=
\begin{pmatrix}0\\0\\0\\1\end{pmatrix}
\quad\Longleftrightarrow\quad \mA\bar\bx=-\bb ,
\end{equation}
where $\mA$ is the upper-left $3\times3$ block of $\bar\mQ$ and $\bb$ the first three entries of its
last column. If $\mA$ is invertible, $\bar\bx=-\mA^{-1}\bb$. If it is singular (the planes do not
determine a point, for instance on flat or cylindrical regions), choose the best among $\bx_1$,
$\bx_2$ and their midpoint, or use the pseudo-inverse. The level sets of $\Delta$ are ellipsoids
whose shape reveals the local geometry: thin along the normal, elongated along flat directions.

\begin{example}{Corners and flat regions}{qem}
At the corner of a cube the vertex has the three planes $x=1$, $y=1$, $z=1$. Then
$\mA=\mI$ and $\bb=-(1,1,1)\T$, so the optimal position is the corner itself with zero error:
collapsing an edge into a corner keeps the corner exactly. For a vertex in the middle of a face
$z=1$, all planes coincide, $\mA=\diag(0,0,1)$ is singular, and any point of the plane has zero
error: flat regions can be decimated freely.
\end{example}

\Cref{alg:qem} summarises the method. A priority queue keyed by the collapse cost gives
$O(n\log n)$ complexity, and each vertex stores only the ten distinct entries of its quadric. Before
performing a collapse, legality checks reject those that would violate the link condition
(\cref{ch:ds}), flip the normal of a neighbouring triangle, or create degenerate triangles. Open
boundaries are preserved by adding, for each boundary edge, a heavily weighted plane perpendicular
to the surface through that edge.

\begin{algorithm}[tbp]
\caption{Simplification with quadric error metrics \citep{Garland1997}.}\label{alg:qem}
\Input{mesh $\M$, target number of faces $F^\star$}
compute $\mQ_v$ for every vertex\;
\ForEach{edge $(v_1,v_2)$}{compute $\bar\bx$ minimising $\bar\bv\T(\mQ_1+\mQ_2)\bar\bv$ and push the edge with cost $\Delta(\bar\bv)$ into a priority queue\;}
\While{$F>F^\star$}{
  pop the edge of smallest cost; skip it if it is stale or the collapse is illegal\;
  collapse it, place the new vertex at $\bar\bx$ and set its quadric to $\mQ_1+\mQ_2$\;
  recompute the cost of the edges incident to the new vertex and update the queue\;
}
\end{algorithm}

\fig{ch10_qem_levels.png}{QEM simplification of the bunny from 20\,000 faces to 2000, 500 and
150. Even at 150 faces the ears and the overall shape survive, because collapses happen first
where the planes agree.}{fig:qem-levels}

\fig{ch10_qem_compare.png}{Comparison of simplification criteria on the cow. Left: mean distance
to the original mesh (per mille of the bounding-box diagonal) against the number of faces. Right:
results with about 200 faces (QEM and shortest edge first exactly 200, vertex clustering 225 with
the closest grid size). QEM is consistently the most accurate; clustering loses the legs and the
head first.}{fig:qem-compare}

\paragraph{Progressive meshes.} Recording the sequence of collapses yields a multiresolution
representation: the coarsest mesh plus the inverse operations (vertex splits), which can be
streamed and replayed to any level of detail \citep{Hoppe1996}.

\section{Subdivision}

Subdivision generates a smooth surface as the limit of a sequence of meshes. Each step
\begin{enumerate}
  \item \textbf{refines} the connectivity, splitting faces (\emph{primal} schemes) or replacing
  vertices by faces (\emph{dual} schemes);
  \item \textbf{smooths} the geometry, computing new positions as fixed affine combinations of
  nearby old positions (the \emph{masks} or rules).
\end{enumerate}
\emph{Approximating} schemes move the old vertices, \emph{interpolating} schemes keep them. The
limit surface is smooth everywhere, while the mesh retains arbitrary topology: subdivision offers
the trade-off between the smoothness of splines and the flexibility of meshes \citep{Zorin2000}.

\subsection{Curves}

\paragraph{Chaikin's algorithm} \citep{Chaikin1974} cuts corners: each edge $(\bp_i,\bp_{i+1})$ is
replaced by the two points
\begin{equation}
\bq_i=\tfrac34\bp_i+\tfrac14\bp_{i+1},\qquad \br_i=\tfrac14\bp_i+\tfrac34\bp_{i+1},
\end{equation}
so an open polygon of $n$ points becomes one of $2(n-1)$ points. The limit is the uniform quadratic
B-spline of the control polygon, which is $C^1$.

\paragraph{Cubic B-splines.} Inserting edge midpoints and moving the old vertices,
\begin{equation}
\bp'_{2i+1}=\tfrac12(\bp_i+\bp_{i+1}),\qquad \bp'_{2i}=\tfrac18(\bp_{i-1}+6\bp_i+\bp_{i+1}),
\end{equation}
converges to the uniform cubic B-spline, which is $C^2$.

\paragraph{The 4-point scheme} \citep{Dyn1987} interpolates: old points are kept and each new point
is
\begin{equation}
\bp'_{2i+1}=\tfrac{9}{16}(\bp_i+\bp_{i+1})-\tfrac1{16}(\bp_{i-1}+\bp_{i+2}),
\end{equation}
with a $C^1$ limit (\cref{fig:curves}).

\fig[0.9\textwidth]{ch10_curves.pdf}{Curve subdivision of a closed control polygon (grey). Chaikin
and the cubic B-spline scheme approximate the polygon and converge to smooth curves inside it;
the 4-point scheme interpolates every control point.}{fig:curves}

\subsection{Surfaces}

\paragraph{Loop subdivision} \citep{Loop1987} works on triangle meshes. Every triangle is split into
four by inserting a vertex on each edge (primal $1\to4$ split). For an interior edge with endpoints
$\bv_0,\bv_1$ and opposite vertices $\bv_2,\bv_3$, the new \emph{edge point} is
\begin{equation}
\be=\tfrac38(\bv_0+\bv_1)+\tfrac18(\bv_2+\bv_3),
\end{equation}
and an old vertex $\bv$ of valence $n$ with neighbours $\bv_j$ moves to the \emph{vertex point}
\begin{equation}
\bv'=(1-n\beta)\bv+\beta\sum_{j=1}^n\bv_j,\qquad
\beta=\frac1n\Bigl(\frac58-\bigl(\tfrac38+\tfrac14\cos\tfrac{2\pi}{n}\bigr)^2\Bigr) .
\end{equation}
For $n=6$, $\beta=1/16$ and the scheme reproduces quartic box splines, which are $C^2$; at
extraordinary vertices ($n\neq6$) the limit is $C^1$. A simpler choice is $\beta=3/16$ for $n=3$ and
$\beta=3/(8n)$ otherwise. Boundary edges and vertices follow the cubic B-spline curve rules, so
boundaries converge to cubic curves.

\paragraph{Catmull--Clark subdivision} \citep{Catmull1978} works on arbitrary polygon meshes and
generalises bicubic B-splines. One step creates
\begin{itemize}
  \item a \emph{face point} for each face: the average of its vertices;
  \item an \emph{edge point} for each edge: the average of its two endpoints and the two face points
  of its adjacent faces;
  \item a new position for each old vertex of valence $n$:
  $\bv'=\dfrac{\mathbf F+2\mathbf R+(n-3)\bv}{n}$, where $\mathbf F$ is the average of the adjacent
  face points and $\mathbf R$ the average of the midpoints of the incident edges.
\end{itemize}
Each face is connected to its edge points and face point, so after one step every face is a
quadrilateral. The limit is $C^2$ except at extraordinary vertices (valence $\neq4$), where it is
$C^1$. Catmull--Clark surfaces are the standard of film and animation modelling.

\paragraph{Doo--Sabin subdivision} \citep{Doo1978} is a dual scheme that generalises biquadratic
B-splines. For each face and each of its vertices a new point is created, for example the average
of the vertex, the midpoints of its two edges in the face and the face centroid. New faces are then
formed of three kinds: a shrunken copy of every face, a quadrilateral for every edge, and an
$n$-gon for every vertex of valence $n$. The limit is $C^1$.

\begin{table}[tbp]
\centering\small
\begin{tabular}{@{}lllll@{}}
\toprule
Scheme & Mesh & Type & Behaviour & Smoothness (regular / extraordinary) \\
\midrule
Chaikin & curves & dual & approximating & $C^1$ \\
Cubic B-spline & curves & primal & approximating & $C^2$ \\
4-point & curves & primal & interpolating & $C^1$ \\
Loop & triangles & primal & approximating & $C^2$ / $C^1$ \\
Catmull--Clark & polygons (quads) & primal & approximating & $C^2$ / $C^1$ \\
Doo--Sabin & polygons & dual & approximating & $C^1$ / $C^1$ \\
Butterfly & triangles & primal & interpolating & $C^1$ (modified) \\
\bottomrule
\end{tabular}
\caption{Common subdivision schemes.}
\label{tab:subdiv}
\end{table}

\fig[0.85\textwidth]{ch10_subdivision.png}{Subdivision of a cube with Catmull--Clark, Doo--Sabin and
Loop (applied to the triangulated cube) at levels 0, 1, 2 and 4. All three converge to rounded
shapes; Loop's limit is less symmetric because the triangulation of the cube is not
symmetric.}{fig:subdivision}

\fig{ch10_loop_bunny.png}{Loop subdivision of a 400-face bunny. Each level multiplies the number of
faces by four; the limit is smooth but, since the scheme approximates, it shrinks slightly away from
the control mesh.}{fig:loop-bunny}

\paragraph{Limit positions and evaluation.} The rules are linear, so the vertices near an
extraordinary vertex are transformed by a fixed \emph{subdivision matrix} at every step. Its
eigen-decomposition gives the limit position (the left eigenvector of the dominant eigenvalue
$1$ applied to the control points) and the tangent plane, and determines the smoothness of the
limit surface; this analysis is how the weights $\beta$ above were chosen.

\begin{keypoints}
\begin{itemize}[leftmargin=*]
  \item Incremental decimation collapses the edge of smallest error first; QEM measures the
  squared distance to the planes of the original neighbourhood with a $4\times4$ quadric
  $\mQ=\sum\bp\bp\T$ and places the new vertex by solving a $3\times3$ system.
  \item Collapses must preserve the topology and avoid normal flips; quadrics add, so updates are
  cheap.
  \item Subdivision alternates refinement and smoothing; Loop (triangles), Catmull--Clark (quads)
  and Doo--Sabin (dual) produce $C^1$ or $C^2$ limit surfaces from coarse control meshes.
\end{itemize}
\end{keypoints}

\section*{Exercises}
\begin{exercise}
Show that the quadric error of a vertex is invariant under the choice of sign of the plane
normals, and that summing quadrics of coplanar faces counts the same plane several times. How
does area weighting help?
\end{exercise}
\begin{exercise}
Verify that Loop's $\beta$ equals $1/16$ for $n=6$ and that the vertex rule then coincides with the
regular quartic box-spline mask.
\end{exercise}
\begin{exercise}
Count vertices, edges and faces after one Catmull--Clark step on a mesh with $V$, $E$, $F$ and check
that the Euler characteristic is preserved.
\end{exercise}
